\answer{ex:07:1}{$V=\kappa p/\bigl(\kappa p+(1-\kappa)(1-p)\bigr)$: $0.059$, $0.20$, $0.50$; la mediana es $0$, mientras que el punto~\eqref{eq:expectile} varía suavemente con $\kappa$.}
\answer{ex:07:2}{$p=1/3$, $x=0.75$, $\sigma_r=0.35$, $V(0.2)=0.083$.}
\answer{ex:07:3}{$V=\sqrt\kappa/(\sqrt\kappa+\sqrt{1-\kappa})$, de $0.25$ a $0.75$; dispersión $\propto\sqrt\eta$; para las gotas $V=0.25+0.5\kappa$; las distribuciones difieren en $0.05$ en las neuronas extremas, hace falta $\eta\lesssim0.01$.}
\answer{ex:07:4}{(a)~$J^*=2\sqrt{C\bar R}$. (b)~Sobrecoste $\bigl(k^{1/2}+k^{-1/2}\bigr)/2-1$: $6\%$ con $k=2$ y $25\%$ con $k=4$; basta una precisión “dentro de un factor dos”.}
\answer{ex:07:5}{Pico de área $R(e^{-1}-e^{-2})\approx0.23R$, como $2.3\,\text{s}$ de rampa; si la dopamina informa de $V$, no hay pico: el nivel simplemente sube en escalón $e$ veces.}
\answer{ex:07:6}{El suceso desplaza el peso $\propto A\tau_f/(\tau_e+\tau_f)$, el nivel da un error $s\tau_f$: $9\,\text{s}\le\tau_f\le17\,\text{s}$.}
\answer{ex:07:7}{De partida $\Delta P=0.80$, $M=0.90$. (a)~$\Delta P=0.74$ ($-8\%$), $M=0.43$ ($-52\%$); (b)~$\Delta P=0.40$ ($-50\%$), $M=0.80$ ($-11\%$): doble disociación.}
\answer{ex:07:8}{$N_{\text{ef}}\to2+\rho^2N\approx10^6$ intentos: $10^4$ veces menos que con un solo cable, pero una fuga del $1\%$ sigue costando un millón de intentos.}
